\section{行为 - 危害行为}
\begin{enumerate}
  \item 3 个特征。
        \begin{enumerate}
          \item 有体性。\yi 必须是\zd{人的行为}活动。\zd{积极举动、消极静止}。
                \al{甲与乙在船上发生争吵，甲追打乙不慎落水，请求施救，乙就干站着}。
                \er 思想不能犯罪。
          \item 有意性。要\zd{有意识}去实施。\al{梦游 \cuo}。
          \item 有害性。\zd{法益侵害性}。刑法保护的界限有特定要求，要和道德区分。\al{兽奸罪、通奸罪}。
        \end{enumerate}
  \item 判断。\bs{对法益是否有危险}。一定要\bs{先客观，后主观}。排除\bs{生活、不能、自陷、升降梯}。
        \begin{enumerate}
          \item 生活行为。
                \begin{enumerate}
                  \item \al{送溜冰鞋，希望溜冰时候摔死}。
                  \item \al{雷雨天，小芳约狗蛋到大树下见面，久等不来被劈死}。
                \end{enumerate}
          \item 升高危险。
          \item 降低危险。\al{广告牌掉下来，狗蛋故意轻轻推小芳}。
          \item vs 替代危险。\al{火灾扔婴儿}。
                \begin{enumerate}
                  \item 都降低了危害结果。前者致命危险 \yjt 重伤危险。
                  \item 但达成方式不同。降低针对\zd{同一危险源}。替代创设了\zd{新的危险源}（火灾 \yjt 高空坠落）。
                \end{enumerate}
          \item 被害人自陷风险。\bs{严格 2 步走}。\zd{危险的实行者、支配者}对风险是否\zd{有认知 + 控制能力} \yjt \zd{推着风险往前}。
                \begin{enumerate}
                  \item \al{你游到河对岸就做你女朋友，不幸溺水死亡}。
                  \item \al{甲将 80g 毒品送给毒瘾很大的朋友乙，乙因过量吸食死亡}。
                  \item \al{为抓捕犯人跑向马路对面}。
                  \item \al{欠钱不还拒之门外，讨债人冻死}。
                  \item \al{女朋友让男朋友闯红灯，结果出车祸女朋友死了}。
                  \item \al{台湾男开车让女口车祸案 \dui 女方有认知，但无控制能力}。
                  \item \bs{放火。不管谁跑进去救人还是拿贵重物品，统统归结于放火行为}。
                \end{enumerate}
          \item 不能犯/不可罚的不能犯。身份/对象不能、手段不能。\al{沙漠杀稻草人}。
        \end{enumerate}
\end{enumerate}
